%*******************************************************
% Abstract in German
%*******************************************************
\begin{otherlanguage}{ngerman}
	\pdfbookmark[0]{Zusammenfassung}{Zusammenfassung}
	\chapter*{Zusammenfassung}
Heutzutage gewinnt "Digitale Souveränität" in der EU immer mehr an politischer und gesellschaftlicher Bedeutung. Cloud Computing ist dabei eine der zentralen Schlüsseltechnologien zur digitalen Transformation, aber es birgt einen grundlegenden Widerspruch: Beim Einsatz von Rechenresourcen in Cloud gibt man die Kontrolle von seinen digitaler Infrastruktur an den externen Cloud-Anbieter ab, was dem Anspruch digitaler Souveränität direkt widerspricht. 

Diese Hausarbeit untersucht die Forschungsfrage: "Wie die digitale Souveränität der EU trotz der Abhängigkeit von US-amerikanischen Cloud-Diensten gewährleistet werden kann?". Auf Basis qualitativer Methoden und systematischer Literaturrecherche werden die wirtschaftlichen, rechtlichen, technischen und geopolitischen Implikationen dieser Abhängigkeiten analysiert sowie bestehende regulatorische Lösungsansätze der EU kritisch bewertet.

Ziel ist es, eine Methode zu finden, die die Vorteile von Cloud-Computing nutzt, ohne die Kontrolle über kritische digitale Infrastrukturen zu verlieren. Die Hausarbeit kommt zu dem Schluss mit dem Konzept der selektiven Souveränität als praktikabelster Ansatz empfohlen: Kritische Infrastrukturen sollten bevorzugt bei europäischen Anbietern oder in hybriden Umgebungen betrieben werden, für sensible Daten wird clientseitige Verschlüsselung empfohlen, und das Jurisdiktionsproblem des US-Cloud-Act sollte durch eine bilaterale Vereinbarung mit den USA adressiert werden. 

\end{otherlanguage}
